\section{\sys: A Question-Driven Kernel Agent}
\label{sec:workflow}
% Question counts, 162 messages, 71 tool calls per message come from paper/analysis/traces.

\sys takes a working CUDA kernel, a set of scenarios, and a goal, and returns a faster kernel whose
outputs match the seed on every scenario. It has three parts: a minimal harness that makes the goal
verifiable, a general-purpose coding agent with no script around it, and a source of questions at the
moments the agent is about to go wrong. The questions come from a human steerer or, once they recur,
from a bank the agent applies to itself. Nothing in the method is specific to an operator, and
nothing requires a faster implementation to exist.

\subsection{Ask vs.\ Equip: A Common Frame}
\label{sec:ask-vs-equip}
\label{sec:frame}

Kernel optimization by an agent is a search. A candidate is a kernel source; its fitness is a vector
of per-scenario speedups over the seed, collapsed to zero if any scenario is incorrect; and the
lineage is monotone, as in AVO \cite{chen2026avo}. The agent is the variation operator: it proposes
the next candidate conditioned on a context, and everything a designer can do is a choice of what
enters that context and when. An \emph{equipped} agent receives the task, tools whose output has
been distilled for it (a profiler parser that names the bottleneck, a disassembly differ), knowledge
retrieved for it (documentation sections, hints mined from reference libraries, a store of lessons),
the lineage, and feedback on the last candidate; the designer improves it by adding tools and
knowledge. Our baseline (\S\ref{sec:baseline}) is a serious instance. An \emph{asked} agent receives
the task, the harness's output, the lineage, and, at decision points, a short question. It still has
every tool a developer machine offers, the compiler, profiler, disassembler, sanitizer, and
debugger, but nothing is pre-digested, and direction arrives as a question rather than a menu. The
designer improves it by asking better questions.

\subsection{A Minimal Harness and a Verifiable Goal}
\label{sec:workflow-harness}

The only component we build is the evaluator. One command builds a candidate, checks its outputs
against the seed's and an fp32 reference on every scenario, and times it against the seed back to
back on the same GPU, reporting per-scenario ratios and their geometric mean. Other kernels the user
cares about can be added as measured comparison columns, but none is required. Constraints are
stated once and enforced mechanically: in the case study, the ban on CUTLASS, CuTe, cuBLAS, and
cuDNN is checked by inspecting the built binary's symbols.

The goal is a predicate over the evaluator's output, enforced rather than requested. We install it
as a Stop hook in the coding agent: whenever the agent attempts to end its turn, the hook evaluates
the predicate and, if any clause is unmet, returns the agent to work with the unmet clauses in hand.
A predicate such as ``at least parity with the seed on every scenario'' removes the most common
failure of autonomous agents, declaring completion early. Two lessons shaped how we write
predicates. They must quantify over every scenario rather than an aggregate, because an agent will
otherwise satisfy the aggregate on the shapes it is good at. And they must not be judged by a model
where they can be computed: the one time our hook delegated the verdict to a model judge, it
accepted a premature claim that a human then overturned (\S\ref{sec:eval-cases}).

\subsection{Questions: What They Are and When to Ask}
\label{sec:workflow-questions}
\label{sec:good-questions}
\label{sec:workflow-when}


Across the traces, the questions that preceded the largest gains share six properties. They carry
no answer: they name a class of hypothesis and leave the derivation to the agent. They demand a
measurement the agent can make with the tools it has, so the answer is a number rather than an
opinion. They move across layers of the system, from the kernel to the host launch path to the
routing logic to the benchmark itself, where an agent left alone stays at the layer it started at.
They refuse premature closure by asking whether a claimed limit is a property of the hardware or of
the agent's own design. They are cheap, tens of tokens against the tens of thousands that retrieved
documentation costs \cite{liu2024lost,langchain2025toolload}. And they counteract deference: an
agent tends to keep the structure it started from and accept the framing it is given
\cite{sharma2024sycophancy}, and a question asks it to justify that choice with evidence.

Eight categories cover the questions a human asked during the case study
(Table~\ref{tab:qbank-human} in Appendix~\ref{app:questions}): measurement check (13 messages), bound
differencing (10), benchmark validity (7), gap triage (7), agent process (6), challenge a limit (5),
strategy options (5), and hypothesis check (3). The largest checks a measurement: whether a win is noise, whether a comparison is fair, whether the
agent is timing what it thinks it is timing. The second pits the agent's claimed ceiling against an
independent estimate of what is achievable. In the case study that estimate was often the upstream
kernel, but the pattern does not depend on one existing: a roofline, a speed-of-light figure
\cite{lin2026solexecbench}, or an A/B that swaps only the launcher around an identical binary serves
the same role. Questions are asked at
decision points, not continuously; the traces show five: at the start of a campaign, to fix the
scope; at a plateau; when the agent claims a limit or completion; after a suspiciously large win;
and when correctness breaks in a way the agent explains away. Between these points the agent works
alone, about 71 tool calls per human message, where the median message is a sentence or two. The
steerer contributes no code and no CUDA knowledge the agent lacks; it contributes doubt, at the right
moment, about the right thing.

\subsection{The Question Bank}
\label{sec:workflow-bank}

After the campaigns we distilled the recurring questions into a bank. Its first version, taken from
the traces alone, had 48 entries, and reading it back showed the hazard of distilling from one case
study: 21 named the upstream kernel as the comparison and 13 named attention structure. The released
bank removes both. A question is admitted only if it carries no answer, is answerable by a
measurement with stock tools, is phrased against a bound (a roofline or a per-tile cycle budget)
rather than a reference, names an architecture only as a parameter, and names no operator. To cover
the aspects a single attention campaign never exercised, we also mined three public GEMM worklogs on
Ampere, Hopper, and Blackwell \cite{boehm2022cudammm,shankhdhar2024h100,modular2025blackwell} and
the FlashAttention-3 and -4 papers \cite{shah2024flashattention3,zadouri2026flashattention4} for
the diagnostic that preceded each of their optimizations, and phrased it as a question. The result
has 101 core entries in ten groups ordered by where the gap might be (scope and triage; the bound
and the regime; harness and measurement; host, launch, and routing; memory hierarchy; compute issue;
occupancy and registers; asynchrony and pipelining; scheduling; low-precision numerics), five that
apply only when another implementation is available to measure, and eighteen per-architecture
checklist entries that bind the generic mechanisms to Ampere, Hopper, or Blackwell instructions.
Operator structure enters through a plug-in with five slots (regimes, irregular work, non-MMA
units, reduction splits, indirect loads); the attention plug-in fills them with prefill and decode,
causal masks, softmax, split-KV, and paged gathers. Appendix~\ref{app:questions} lists the bank with
provenance tags: 29 entries carry a trace tag, and the 72 from public sources alone are untested.

\subsection{From Asking the Agent to the Agent Asking Itself}
\label{sec:workflow-skill}

Because the questions recur, they can be handed to the agent as part of its instructions. We codify
them as a skill with five parts: the goal, the hard rules the harness enforces, a loop, a list of
questions to ask itself, and the commands that run the evaluator and profiler. The loop is the bank
in procedural form: fix the scope; measure the seed; classify each gap as host, routing, or kernel
before touching code; bound it; form a hypothesis; rewrite; gate. At each step the agent fans out
subagents, one per hypothesis, each returning a hypothesis and the measurement that tests it rather
than a patch. The skill is parameterized by the seed, the scenario set, and the goal; its questions
were reused across the Hopper and Blackwell campaigns nearly verbatim.

\subsection{Composing with Evolutionary Search}
\label{sec:workflow-evo}

\sys is compatible with the evolutionary loops of the open-ended agents, because it changes the
variation operator's direction rather than the loop around it. We wrap the same agent in an
AVO-style search \cite{chen2026avo}: a monotone lineage, several diverse mutations per generation
applied to isolated copies and scored by the evaluator, and a strategist that reviews the trajectory
and writes the next generation's directives. The one change is that the directives are questions,
hypotheses to test rather than edits to make. As \S\ref{sec:eval-evo} shows, evolution amplifies direction; it does not
supply it.

\subsection{Hypotheses}
\label{sec:hypotheses}

The evaluation tests five claims. \textbf{H1 (quality):} on the same task, evaluator, and hardware,
the asked agent reaches higher performance at full correctness than the equipped one.
\textbf{H2 (efficiency):} with fewer model-driven steps and GPU-hours, counting human time.
\textbf{H3 (mechanism):} the largest gains follow questions that change the layer under
consideration or reject a claimed limit, not questions that supply information.
\textbf{H4 (complementarity):} evolutionary search amplifies the effect of questions and does not
substitute for them. \textbf{H5 (generality):} the questions that moved the case study contain
nothing specific to the operator or to a reference once the reference is replaced by a bound.
